\chapter{Conclusion and Recommendations}
\label{chap:conclusion}

For this study, a computational model based on the DPWH 2024 Atlas was created to detect and quantify the existence of Braess's Paradox within Iloilo City. This provided a new methodology of evaluating the traffic management system of a complex traffic network in an urban Iloilo City.

To construct this sytem, primary arterial roads and intersections were identified as nodes and edges, respectively. The volume and direction of travel demand was applied to the network from the OD matrix pairs by zonal areas, with the Annual Average Daily Traffic (AADT) served as weights amongst the edges.

The Price of Anarchy (PoA) was calculated to determine the critical volume for the identification of Braess's Paradox in the network, while the Bureau of Public Roads (BPR) function was used to update the travel of each edge. The System Optimum baseline (SO) was established to apply systematic delay from the vehicles into the network. 

The study underwent several experimental setups using the computational model. The Zonal OD pairs were implemented in the network by using centroid nodes. There are three connector topologies used in the study: uniform spatial disaggregation, unrestricted connector, and AADT-prioritized. These connector topologies were then tested with two methods of road closure: bidirectional and unidirectional.

As a result, the computation model shows that the results reflect the existence of Braess Paradox links within the Iloilo City network. In the Price of Anarchy (PoA) recorded at 1.581, the results show that selfish routing behavior can lead to a 58.15\% increase in the Total System Travel Time (TSTT) compared to a socially optimal state. The results also identified 4,627 unidirectional link closures and 3,662 bidirectional link closures that mathematically improved the Total System Travel Time (TSTT) within the model. Amongst the six topologies tested, the AADT-prioritized topology with unidirectional closure has the most identified of 21 Braess Paradox links.

Additionally, it was found that the traffic congestion in Iloilo City is due to the directional imbalances and unnecessary roads, which leads to drivers going to inefficient routes. The result of the unidirectional closures suggests that restricting traffic to one direction can lead to significant time savings as it reduces conflicts in going to major intersections. The bidirectional closures of roads have removed redundant paths that hinder the driver from taking the selfish routes and lead the traffic into a more efficient network. 

In the evaluation of the six topologies, the volume of traffic is diluted and no paradoxes occur when traffic is spread out evenly in the city. When a demand is placed through high-volume intersections for realistic peak hours, paradoxes appear that prove that the congestion in Iloilo City is volume-dependent through directional sensitivity rather than a lack of road space. The result in the topology also shows that unidirectional closures are often more efficient than full bidirectional closures. This is due to how a restriction of a single direction can resolve intersection conflicts without causing widespread traffic spills, unlike with total road closures.

Based on the findings of the study, future research should consider different networks that have accessible and available traffic data. Utilizing other networks allows comparison of accuracy of the traffic network to the reflected traffic data. High-quality data, instead of simplified ones like zonal OD matrix, would also improve the reliability of the computational model better suited for the methodology framework used by Youn et al. (2008).

Vehicles in a traffic network do not only traverse through the primary arterial roads. Further improvement of the computational modeling is to implement secondary and tertiary roads, as secondary roads will act as shortcuts. This recommendation oversees conditions where high-capacity primary roads are avoided in favor of lower-capacity secondary roads that triggers a bottleneck, therefore presenting the paradox. 

Public transport systems also occupy the traffic network while also contributing to its complexity by their unpredictable operational behavior. The inclusion of public transport in further research can help in creating a more realistic traffic network model while also considering it as a variable that may contribute to further appearance of the Braess paradox link. Since public vehicles follow fixed routes, future studies may implement them as a fixed flow that may reduce the width of the road for the private vehicles. 

In any traffic network, traffic volume naturally fluctuates between  heavy flows during rush hours and lighter ones during off-peak hours. While the study simulates the local traffic flow in Iloilo City, it does not account for the seasonality of the traffic data. Future studies should collect and incorporate traffic data that is categorized by various times of the day, similarly between peak and non-peak hours. By comparing these traffic windows, this would provide more insights on how congestion moves throughout the day in the network. Additionally, the data would give a more accurate representation of the traffic network in Iloilo City.

To address the computational limitations faced in this study, future research should explore alternative libraries and software tools. Furthermore, researchers may create their own custom data structures to optmize processing speed and memory management when simulating large-scale traffic networks.

Finally, discussion with stakeholders about the expansion of this study is highly recommended. The stakeholders such as the local government would help in creating a model that caters to the traffic network of Iloilo City rather than relying on standard coefficients such as the $\alpha$ and $\beta$ parameters in the BPR function. Presenting the findings to the Iloilo City Traffic and Transportation Management may help in evaluating proposed road projects.